\documentclass[11pt]{article}
\usepackage[margin=1in]{geometry}
\usepackage{amsmath,amssymb,graphicx,booktabs,hyperref,xcolor}
\title{Computational Replication Report:\\ Ding \emph{et al.} (2021/2022),\\ \emph{Observation of the Orbital Rashba-Edelstein Magnetoresistance} (arXiv:2105.04495, PRL 128)}
\author{Ollie (OpenClaw @ Argonne) --- TEXTURES-100 Wave 4}
\date{2026-07-19}
\begin{document}
\maketitle

\section{Paper summary}
The paper reports magnetoresistance (MR) from orbital-angular-momentum (OAM) transport --- the
\emph{orbital Rashba-Edelstein magnetoresistance} (OREMR) --- in a Permalloy/oxidized-Cu
(Py/Cu$^*$) heterostructure, notably \emph{without} heavy (strong-SOC) elements. The decisive
evidence is the angular dependence of the MR in three orthogonal field-rotation planes
($\alpha=xy$, $\beta=yz$, $\gamma=zx$): ordinary AMR gives a flat $\beta$ scan, but the observed
$\rho_{xx}(\beta)\sim\cos^2\beta$ is SMR-like and signals an interfacial orbital Rashba-Edelstein
effect. The MR ratio is thickness-independent below the oxidation depth (interfacial signature).

\section{Scope statement}
This is an experimental paper. We replicate the \emph{theory framework}: the SMR/OREMR angular
resistivity model, the three-plane $\cos^2$ laws, the $\beta$-scan discriminator that identifies
OREMR over AMR, and the interfacial thickness signature. Measured MR magnitudes and the
spin-diffusion/dephasing-length extraction require the physical samples (out of scope).

\section{Model}
Longitudinal SMR/OREMR resistivity (current $\|x$, Chen-type form):
\[
\rho_{xx}(\mathbf m)=\rho_0 + \Delta_{\rm AMR}\,m_x^2 + \Delta_{\rm SMR}\,(1-m_y^2),
\]
with $\Delta_{\rm AMR}{=}0.02$, $\Delta_{\rm SMR}{=}0.015$ (illustrative). AMR-only sets
$\Delta_{\rm SMR}{=}0$.

\section{Claims addressed}
\begin{table}[h]\centering\small
\begin{tabular}{clccc}
\toprule
ID & Claim & Type & Testable? & Tested? \\
\midrule
C1 & $\beta$ scan flat (AMR) vs $\cos^2\beta$ (OREMR): discriminator & theory & yes & yes \\
C2 & Three-plane $\cos^2$ angular laws, high $R^2$ & theory & yes & yes \\
C3 & MR ratio interfacial (flat for $t_{\rm Cu^*}\!\le\!5$nm) & theory/model & yes & yes \\
C4 & Measured MR magnitudes; diffusion/dephasing lengths & experiment & no & \textcolor{red}{out of scope} \\
\bottomrule
\end{tabular}
\end{table}

\section{Results vs.\ paper}
\begin{table}[h]\centering\small
\begin{tabular}{lccc}
\toprule
Quantity & Expected & This work & Verdict \\
\midrule
$\beta$ scan peak-to-peak (AMR only) & $0$ (flat) & $0.0$ & \checkmark \\
$\beta$ scan peak-to-peak (OREMR) & $>0$ ($\cos^2$) & $0.015$ & \checkmark \\
$\alpha$ scan $\cos^2$ fit $R^2$ & $\to 1$ & $1.000$ & \checkmark \\
$\beta$ scan $\cos^2$ fit $R^2$ & $\to 1$ & $1.000$ & \checkmark \\
$\gamma$ scan $\cos^2$ fit $R^2$ & $\to 1$ & $1.000$ & \checkmark \\
MR-ratio plateau flatness ($t\!\le\!5$nm) & $\approx 0$ & $0.0$ & \checkmark \\
\bottomrule
\end{tabular}
\end{table}

\section{Per-claim what worked / what didn't}
\textbf{C1 (worked).} AMR-only $\beta$ scan is exactly flat; adding the SMR/OREMR term produces the
$\cos^2\beta$ modulation --- the paper's decisive discriminator. \textbf{C2 (worked).} All three
scans fit $\cos^2$ to $R^2=1$. \textbf{C3 (worked).} The MR ratio is flat below the oxidation depth
(interfacial) and decays beyond. \textbf{C4 (out of scope).} Absolute MR ratios and length
extractions need the measured samples.

\section{Critique of the replication}
Honest: C1--C3 are analytic consequences of a phenomenological model with hand-chosen coefficients.
They constitute a \emph{theoretical consistency check} that the OREMR angular-dependence framework
behaves as the paper describes (and correctly distinguishes OREMR from AMR), NOT an independent
reproduction of any measured number. No experimental observable is regenerated. The interfacial
plateau flatness is constructed. The LLM-judge scored PARTIAL (coverage 7, agreement 6), which we
adopt.

\section{Verdict}
\textbf{PARTIAL} --- the SMR/OREMR angular-dependence theory, the $\beta$-scan discriminator, and the
interfacial thickness signature are reproduced; measured MR magnitudes and diffusion/dephasing-length
extractions are experimental and out of scope.

\section*{Open Questions}
See \texttt{report/open\_questions.json}: Q1 orbital vs spin mixing conductance; Q2 orbital dephasing
length; Q3 OREMR/SMR/AMR separation protocol; Q4 shunting vs intrinsic OAM quenching; Q5 generality
across light-metal-oxide interfaces.

\end{document}
